\chapter{Conditional Antiunitary Bridge}

This appendix isolates the antiunitary proposal because it is a conditional subhypothesis, not the foundation of the physical--experiential framework. It should remain dormant unless phenomenological relations are independently shown to require a complex projective representation. The existence of quantum-like models in cognition is not by itself sufficient to justify such a representation of experience; Hilbert-space probability can provide useful contextual models without implying quantum physical consciousness or a literally complex phenomenal state space \cite{pothos2022QuantumCognition}.

Suppose independent encodings are nevertheless established,

\[
\Phi_P:P_{\mathrm{phys}}\to\mathcal H_P,
\qquad
\Phi_E:E_{\mathrm{exp}}\to\mathcal H_E.
\]

A candidate antiunitary relation may be written

\[
\Phi_E(E_{\mathrm{exp}})\simeq\Gamma J\Phi_P(P_{\mathrm{phys}}),
\qquad
J=UK,
\]

where $U$ is unitary and $K$ complex conjugation. Antiunitarity implies

\[
\langle J\psi,J\phi\rangle=\overline{\langle\psi,\phi\rangle}.
\]

Pairwise transition magnitudes are therefore invariant:

\[
|\langle J\psi,J\phi\rangle|^2=|\langle\psi,\phi\rangle|^2.
\]

Pairwise magnitudes cannot distinguish a unitary bridge from an antiunitary one. A three-state Bargmann invariant can. Define

\[
B_{123}=\langle\psi_1,\psi_2\rangle\langle\psi_2,\psi_3\rangle\langle\psi_3,\psi_1\rangle.
\]

A unitary transformation preserves $B_{123}$, whereas an antiunitary transformation sends it to its complex conjugate:

\[
B_{123}\mapsto\overline{B_{123}}.
\]

Hence

\[
\arg B_{123}\mapsto-\arg B_{123}.
\]

This provides a sharp discriminator if, and only if, complex projective phenomenological geometry has been justified before the test. A non-circular evidential sequence would therefore require classical probabilistic models to fail out of sample, a Hilbert-space representation to show robust predictive advantage, genuinely complex phase information to be necessary rather than a real embedding with redundant coordinates, and only then unitary versus antiunitary transformation laws to be compared.

Failure of the antiunitary test would reject this subhypothesis rather than relational dual-aspect realism as a whole. Conversely, success would establish a structural correspondence requiring explanation, not the metaphysical identity of conjugation with consciousness. The exact wave result in Chapter 7 remains what it is: conjugation preserves selected form invariants and reverses current. This appendix asks whether an independently established experiential geometry ever exhibits the same transformation signature.
